% DEFINED:   sec:cnn-bias, sec:pure-python-boundary
% RESOLVED:  sec:inductive-biases (backward into Part I)
% FORWARD (prose-only): Chapter 4 (sequences), Chapter 6 (the Transformer)

\section{Inductive bias, named}\label{sec:cnn-bias}

The convolutional network commits to two priors. Made explicit:
\begin{itemize}
\item \textbf{Locality.} A unit reads a $k \times k$ window, not the whole
  image. Long-range interactions emerge only by stacking conv layers, each
  output cell's receptive field growing with depth.
\item \textbf{Translation equivariance.} The same weights apply at every
  position. Shift the input, shift the output.
\end{itemize}
What it gives up is equally concrete:
\begin{itemize}
\item \textbf{Free interaction between distant pixels.} An MLP can wire the
  upper-left corner directly to the lower-right with a single weight. One conv
  layer cannot represent that at all; deeper stacks can, but only indirectly,
  through receptive-field growth.
\item \textbf{Position-specific filters.} If the absolute location of the
  digit carries information, the convolution has to compensate elsewhere, or
  be given extra channels, or be helped by centering the input before it
  arrives.
\end{itemize}

This is the trade, and it is the same shape of trade \Cref{sec:inductive-biases}
described. An MLP is the most generic supervised model: every input may
interact with every other on equal footing. A CNN asserts that the useful
interactions on an image are local and translation-equivariant, and bakes that
assertion into the weights through structural reuse. Whether the assertion
pays is empirical, not a matter of taste. The CNN matches the MLP with fewer
parameters here, where the spatial structure was already half-extracted by
preprocessing, and it pulls clearly ahead on data where that structure is
intact. On genuinely tabular features, where there is no spatial adjacency to
exploit, the same prior is a lie and the CNN loses. The right prior is the one
that matches the data, and the permuted-pixel control is how we checked that
it does.

This chapter took the architecture axis on its first non-trivial commitment,
to a specific group of symmetries, the $2$-D translations. The next chapter
makes the same kind of move on a different data type. Sequences, whether text
or time series, carry their own notion of locality, the recent past mattering
more than the distant past, and their own translation symmetry, a shift in
time rather than in space. A fixed-context language model, and after it a
recurrent network, will commit to those priors by reusing weights across time
steps instead of across spatial positions. Different data, the same lesson:
structural weight sharing is how a network commits to an invariance, and
committing to the right invariance is how it learns more from less.

\section{What stayed pure Python, and what will not}\label{sec:pure-python-boundary}

The five-nested-loop \texttt{Conv2D} worked here because the problem is small:
an $8 \times 8$ single-channel input, a $3 \times 3$ kernel, four output
channels. The forward pass is about 1300 multiply-adds per example and the
backward pass roughly the same, so 3823 examples over 40 epochs finish in a
few minutes of interpreted Python.

This is the last large layer this book writes in pure Python. For $28 \times
28$ MNIST with several conv layers and tens of channels each, the inner loop
count grows by two orders of magnitude, and the five nested loops stop being
merely unflattering and become intractable. At that point a vectorized
implementation, first in NumPy and eventually with batched tensor operations
in a framework, is genuine relief rather than premature optimization. The math
does not change and the indices do not change; only the implementation moves
from interpreted loops to compiled tensor calls. The chapter on the Transformer
will reach for NumPy for exactly this reason, and will say so plainly when it
arrives. Pure Python has carried the whole library to here; it does not have to
carry it everywhere.

One natural extension is worth a line, in the spirit of \Cref{ch:foundations}'s
closing note on automatic differentiation. The convolution produces feature
maps whose positions track the input, which is equivariance. A classifier
usually wants the opposite, invariance: the digit is a seven wherever it sits.
Averaging each feature map down to a single number, a global average pool,
collapses an equivariant stack into a translation-invariant reader, discarding
where the activation was while keeping that it occurred. The library includes
such a layer; stacking it after the convolution is the natural next step toward
a deeper invariant classifier, and it needs no new machinery, only the
\texttt{Layer} contract already in hand.
